\documentclass[11pt]{article}
\usepackage[a4paper,margin=1in]{geometry}
\usepackage{amsmath,amssymb,amsthm}
\usepackage{mathtools}

\title{Explaining the STAR Case in the Strips Proof for Version 2 of \texttt{shifts}}
\author{}
\date{}

\begin{document}
\maketitle

\section*{Goal}

We want to explain the STAR case of the correctness proof for version 2 of \texttt{shifts}. The statement being proved is:

\[
P(ms,r) : \operatorname{set}(\operatorname{shifts}\ ms\ r) = \operatorname{Strips}(L(r))(\operatorname{set}\ ms).
\]

In the STAR case, this becomes:

\[
\operatorname{set}(\operatorname{shifts}\ ms\ r^*) = \operatorname{Strips}(L(r^*))(\operatorname{set}\ ms).
\]

The point of this note is to explain carefully:

\begin{itemize}
\item what the symbols mean,
\item why the two STAR lemmas are true,
\item what the key semantic identity means,
\item and how the STAR proof uses those lemmas.
\end{itemize}

\section*{Basic notation}

\subsection*{Sets and symbols}

If we write
\[
X \subseteq Y,
\]
this means:

\begin{quote}
every element of \(X\) is also an element of \(Y\).
\end{quote}

So \(X\) is contained in \(Y\).

If we write
\[
X = Y,
\]
this means:

\begin{quote}
\(X\) and \(Y\) have exactly the same elements.
\end{quote}

A common way to prove \(X=Y\) is to prove both
\[
X \subseteq Y
\quad\text{and}\quad
Y \subseteq X.
\]

If we write
\[
X - Y,
\]
this means the \emph{set difference}: all elements that are in \(X\) but not in \(Y\).

If we write
\[
\emptyset,
\]
this means the empty set, the set with no elements.

\subsection*{Why \(X-Y=\emptyset\) means the same as \(X \subseteq Y\)}

This is standard set theory.

The statement
\[
X-Y=\emptyset
\]
means:

\begin{quote}
there is no element of \(X\) that lies outside \(Y\).
\end{quote}

But that is exactly the same as saying:

\begin{quote}
every element of \(X\) lies inside \(Y\).
\end{quote}

And that is exactly the meaning of
\[
X \subseteq Y.
\]

So:
\[
X-Y=\emptyset
\quad\Longleftrightarrow\quad
X \subseteq Y.
\]

In our proof, this is used in the form
\[
\operatorname{Strips}(A,B)-B=\emptyset
\quad\Longleftrightarrow\quad
\operatorname{Strips}(A,B)\subseteq B.
\]

\section*{What \(\operatorname{Strips}(A,B)\) means}

The new \(\operatorname{Strips}\) is defined from
\[
\operatorname{Strip}(A,s)=\{\,s_2 \mid \exists s_1\in A.\ s=s_1@s_2\,\},
\]
and then
\[
\operatorname{Strips}(A,B)=\bigcup\bigl(\operatorname{Strip}(A)\,'\,B\bigr).
\]

The important point is that the removed prefix \(s_1\) is allowed to be empty. So if \([] \in A\), then a string in \(B\) may remain unchanged.

Intuitively:

\begin{quote}
\(\operatorname{Strips}(A,B)\) means: all strings you can get by removing one prefix from \(A\) from strings in \(B\).
\end{quote}

If \(A=L(r)\), then:

\begin{itemize}
\item \(\operatorname{Strips}(L(r),B)\) means one pass through \(r\),
\item \(\operatorname{Strips}(L(r)^*,B)\) means zero or more passes through \(r\).
\end{itemize}

\section*{Operational STAR case}

Version 2 of \texttt{shifts} defines STAR as:
\[
\operatorname{shifts}\ ms\ (r^*)
=
\text{let } ms^\# = (\operatorname{shifts}\ ms\ r)-ms
\text{ in if } ms^\# = [] \text{ then } ms \text{ else } ms @ \operatorname{shifts}\ ms^\#\ (r^*).
\]

This means:

\begin{enumerate}
\item shift once through \(r\),
\item remove the marks that did not really change,
\item call the genuinely new shorter marks \(ms^\#\),
\item if no such marks exist, stop and return \(ms\),
\item otherwise keep \(ms\) and recurse on \(ms^\#\).
\end{enumerate}

So the proof naturally splits into two cases:

\begin{itemize}
\item \( (\operatorname{shifts}\ ms\ r)-ms = [] \),
\item \( (\operatorname{shifts}\ ms\ r)-ms \neq [] \).
\end{itemize}

\section*{Lemma 1}

\subsection*{Statement}

If
\[
\operatorname{Strips}(A,B)-B=\emptyset,
\]
then
\[
\operatorname{Strips}(A^*,B)=B.
\]

Since
\[
\operatorname{Strips}(A,B)-B=\emptyset
\quad\Longleftrightarrow\quad
\operatorname{Strips}(A,B)\subseteq B,
\]
we can also read this lemma as:

\begin{quote}
if one \(A\)-step never leaves \(B\), then any finite number of \(A\)-steps never leaves \(B\), and because star also allows zero steps, the result is exactly \(B\).
\end{quote}

\subsection*{Why this lemma is true}

We prove first:

\[
\text{if } \operatorname{Strips}(A,B)\subseteq B,\ \text{then for all }n,\ \operatorname{Strips}(A^n,B)\subseteq B.
\]

\paragraph{Base case \(n=0\).}
Since \(A^0=\{[]\}\), stripping by \(A^0\) changes nothing, so
\[
\operatorname{Strips}(A^0,B)=B.
\]
Hence
\[
\operatorname{Strips}(A^0,B)\subseteq B.
\]

\paragraph{Induction step.}
Assume
\[
\operatorname{Strips}(A^n,B)\subseteq B.
\]
We want
\[
\operatorname{Strips}(A^{n+1},B)\subseteq B.
\]

Using the concatenation property of \(\operatorname{Strips}\),
\[
\operatorname{Strips}(A^{n+1},B)
=
\operatorname{Strips}(A,\operatorname{Strips}(A^n,B)).
\]

By the induction hypothesis,
\[
\operatorname{Strips}(A^n,B)\subseteq B.
\]

By monotonicity of \(\operatorname{Strips}\) in its second argument,
\[
X\subseteq Y \Longrightarrow \operatorname{Strips}(A,X)\subseteq \operatorname{Strips}(A,Y),
\]
we get
\[
\operatorname{Strips}(A,\operatorname{Strips}(A^n,B))
\subseteq
\operatorname{Strips}(A,B).
\]

Now use the assumption
\[
\operatorname{Strips}(A,B)\subseteq B.
\]

So
\[
\operatorname{Strips}(A^{n+1},B)\subseteq B.
\]

Thus the claim holds for all \(n\).

\paragraph{From powers to star.}
Now use
\[
A^*=\bigcup_{n\ge 0} A^n.
\]
Therefore
\[
\operatorname{Strips}(A^*,B)
=
\operatorname{Strips}\Bigl(\bigcup_{n\ge 0} A^n,\ B\Bigr)
=
\bigcup_{n\ge 0}\operatorname{Strips}(A^n,B).
\]

Each set \(\operatorname{Strips}(A^n,B)\) is contained in \(B\), so their union is also contained in \(B\). Hence
\[
\operatorname{Strips}(A^*,B)\subseteq B.
\]

On the other hand,
\[
B\subseteq \operatorname{Strips}(A^*,B),
\]
because \([]\in A^*\), so stripping the empty string is always allowed and leaves every element of \(B\) unchanged.

So we have both inclusions:
\[
\operatorname{Strips}(A^*,B)\subseteq B
\quad\text{and}\quad
B\subseteq \operatorname{Strips}(A^*,B).
\]

Hence
\[
\operatorname{Strips}(A^*,B)=B.
\]

\subsection*{Meaning of Lemma 1}

Lemma 1 says:

\begin{quote}
if one pass through \(A\) never produces anything outside \(B\), then repeated passes through \(A\) also never produce anything outside \(B\); since star also includes zero passes, star gives exactly \(B\).
\end{quote}

This is exactly what is needed in the easy STAR case.

\section*{Lemma 2: the key semantic identity}

\subsection*{Statement}

The second lemma is the identity
\[
\operatorname{Strips}(A^*,B)
=
B \cup \operatorname{Strips}(A^*)\bigl(\operatorname{Strips}(A,B)-B\bigr).
\tag{$\dagger$}
\]

\subsection*{What ``semantic identity'' means}

Here ``semantic identity'' means:

\begin{quote}
a mathematically true equality about the semantic operator \(\operatorname{Strips}\).
\end{quote}

It is called semantic because \(\operatorname{Strips}\) describes what strings remain after matching prefixes from a language. So it talks about the meaning of the regex, not about the syntax of the program.

This identity is not just guessed. It is a set-theoretic property that the proof establishes using lemmas about powers \(A^n\), monotonicity, and the decomposition of \(A^*\).

\subsection*{What the identity means intuitively}

The left-hand side
\[
\operatorname{Strips}(A^*,B)
\]
means all results obtainable from \(B\) by doing zero or more \(A\)-steps.

The right-hand side says that every such result is of one of two kinds:

\begin{enumerate}
\item it is already in \(B\), corresponding to taking zero \(A\)-steps;
\item or it comes from making one genuinely new \(A\)-step and then continuing with zero or more further \(A\)-steps.
\end{enumerate}

Why do we write
\[
\operatorname{Strips}(A,B)-B
\]
instead of just \(\operatorname{Strips}(A,B)\)?

Because in the new version, empty stripping is allowed. So after one \(A\)-step, some results may still already be in \(B\). The term
\[
\operatorname{Strips}(A,B)-B
\]
keeps only the genuinely new results, the ones that really moved forward.

This exactly matches the operational definition
\[
ms^\#=(\operatorname{shifts}\ ms\ r)-ms.
\]

\subsection*{How the identity is proved}

To prove
\[
\operatorname{Strips}(A^*,B)
=
B \cup \operatorname{Strips}(A^*)\bigl(\operatorname{Strips}(A,B)-B\bigr),
\]
we prove both inclusions.

\subsubsection*{Easy direction}

First show
\[
B \cup \operatorname{Strips}(A^*)\bigl(\operatorname{Strips}(A,B)-B\bigr)
\subseteq
\operatorname{Strips}(A^*,B).
\]

This has two parts.

\paragraph{Part 1.}
\[
B\subseteq \operatorname{Strips}(A^*,B),
\]
because \([]\in A^*\).

\paragraph{Part 2.}
\[
\operatorname{Strips}(A^*)\bigl(\operatorname{Strips}(A,B)-B\bigr)
\subseteq
\operatorname{Strips}(A^*,B).
\]

Why? Because
\[
\operatorname{Strips}(A,B)-B \subseteq \operatorname{Strips}(A,B),
\]
and anything obtained by one \(A\)-step from \(B\), followed by zero or more extra \(A\)-steps, is certainly obtainable by zero or more total \(A\)-steps from \(B\).

So the easy inclusion holds.

\subsubsection*{Hard direction}

Now show
\[
\operatorname{Strips}(A^*,B)
\subseteq
B \cup \operatorname{Strips}(A^*)\bigl(\operatorname{Strips}(A,B)-B\bigr).
\]

The file proves this first for powers:
\[
\operatorname{Strips}(A^n,B)
\subseteq
B \cup \operatorname{Strips}(A^*)\bigl(\operatorname{Strips}(A,B)-B\bigr)
\quad\text{for all }n.
\]

\paragraph{Base case \(n=0\).}
Since
\[
\operatorname{Strips}(A^0,B)=B,
\]
the result is immediate.

\paragraph{Induction step.}
Assume
\[
\operatorname{Strips}(A^n,B)
\subseteq
B \cup \operatorname{Strips}(A^*)\bigl(\operatorname{Strips}(A,B)-B\bigr).
\]

Then
\[
\operatorname{Strips}(A^{n+1},B)
=
\operatorname{Strips}(A,\operatorname{Strips}(A^n,B)).
\]

By monotonicity,
\[
\operatorname{Strips}(A,\operatorname{Strips}(A^n,B))
\subseteq
\operatorname{Strips}\Bigl(A,\ B \cup \operatorname{Strips}(A^*)\bigl(\operatorname{Strips}(A,B)-B\bigr)\Bigr).
\]

Distribute over union:
\[
=
\operatorname{Strips}(A,B)
\cup
\operatorname{Strips}\Bigl(A,\operatorname{Strips}(A^*)\bigl(\operatorname{Strips}(A,B)-B\bigr)\Bigr).
\]

Now use the auxiliary fact
\[
\operatorname{Strips}(A,\operatorname{Strips}(A^*,C))
\subseteq
\operatorname{Strips}(A^*,C).
\tag{$*$}
\]

So
\[
\operatorname{Strips}\Bigl(A,\operatorname{Strips}(A^*)\bigl(\operatorname{Strips}(A,B)-B\bigr)\Bigr)
\subseteq
\operatorname{Strips}(A^*)\bigl(\operatorname{Strips}(A,B)-B\bigr).
\]

Hence
\[
\operatorname{Strips}(A^{n+1},B)
\subseteq
\operatorname{Strips}(A,B)
\cup
\operatorname{Strips}(A^*)\bigl(\operatorname{Strips}(A,B)-B\bigr).
\]

Finally, every element of \(\operatorname{Strips}(A,B)\) is either already in \(B\), or lies in \(\operatorname{Strips}(A,B)-B\). So
\[
\operatorname{Strips}(A,B)
\subseteq
B \cup (\operatorname{Strips}(A,B)-B).
\]

And since \(A^*\) contains the empty word,
\[
\operatorname{Strips}(A,B)-B
\subseteq
\operatorname{Strips}(A^*)\bigl(\operatorname{Strips}(A,B)-B\bigr).
\]

So
\[
\operatorname{Strips}(A^{n+1},B)
\subseteq
B \cup \operatorname{Strips}(A^*)\bigl(\operatorname{Strips}(A,B)-B\bigr).
\]

This finishes the induction on \(n\). Since
\[
A^*=\bigcup_{n\ge 0} A^n,
\]
the desired inclusion for \(A^*\) follows.

Thus the identity \((\dagger)\) is proved.

\subsection*{Meaning of the auxiliary fact \((*)\)}

The fact
\[
\operatorname{Strips}(A,\operatorname{Strips}(A^*,C))
\subseteq
\operatorname{Strips}(A^*,C)
\]
means:

\begin{quote}
if you first do zero or more \(A\)-steps, and then do one more \(A\)-step, you are still within the results of doing zero or more \(A\)-steps.
\end{quote}

This is because
\[
A^*=\{[]\}\cup A^*@A,
\]
so one more \(A\) after an \(A^*\)-computation is still part of an \(A^*\)-computation.

\section*{Using the lemmas in the STAR proof}

Now let
\[
A=L(r), \qquad B=\operatorname{set}\ ms.
\]

We want to prove
\[
\operatorname{set}(\operatorname{shifts}\ ms\ r^*)=\operatorname{Strips}(L(r^*))(\operatorname{set}\ ms).
\]

\subsection*{Case 1: \((\operatorname{shifts}\ ms\ r)-ms=[]\)}

Operationally, version 2 of STAR returns
\[
\operatorname{shifts}\ ms\ r^* = ms.
\]
So the left-hand side is
\[
\operatorname{set}(\operatorname{shifts}\ ms\ r^*)=\operatorname{set}\ ms = B.
\]

By the induction hypothesis for \(r\),
\[
\operatorname{set}(\operatorname{shifts}\ ms\ r)=\operatorname{Strips}(L(r),B)=\operatorname{Strips}(A,B).
\]

The assumption gives
\[
\operatorname{Strips}(A,B)-B=\emptyset.
\]

By Lemma 1,
\[
\operatorname{Strips}(A^*,B)=B.
\]

Since \(A=L(r)\), this gives
\[
\operatorname{Strips}(L(r^*),\operatorname{set}\ ms)=\operatorname{set}\ ms.
\]

So the easy STAR case is proved.

\subsection*{Case 2: \((\operatorname{shifts}\ ms\ r)-ms\neq[]\)}

Operationally,
\[
\operatorname{set}(\operatorname{shifts}\ ms\ r^*)
=
B \cup \operatorname{set}(\operatorname{shifts}((\operatorname{shifts}\ ms\ r)-ms,\ r^*)).
\]

By the induction hypothesis for the recursive STAR call,
\[
\operatorname{set}(\operatorname{shifts}((\operatorname{shifts}\ ms\ r)-ms,\ r^*))
=
\operatorname{Strips}(A^*)\bigl(\operatorname{set}((\operatorname{shifts}\ ms\ r)-ms)\bigr).
\]

By the induction hypothesis for \(r\),
\[
\operatorname{set}(\operatorname{shifts}\ ms\ r)=\operatorname{Strips}(A,B).
\]

So
\[
\operatorname{set}(\operatorname{shifts}\ ms\ r^*)
=
B \cup \operatorname{Strips}(A^*)\bigl(\operatorname{Strips}(A,B)-B\bigr).
\]

Now apply Lemma 2:
\[
\operatorname{Strips}(A^*,B)
=
B \cup \operatorname{Strips}(A^*)\bigl(\operatorname{Strips}(A,B)-B\bigr).
\]

Therefore
\[
\operatorname{set}(\operatorname{shifts}\ ms\ r^*)
=
\operatorname{Strips}(A^*,B)
=
\operatorname{Strips}(L(r^*))(\operatorname{set}\ ms).
\]

So the second STAR case is proved.

\section*{Final intuition}

The two lemmas say exactly this:

\begin{itemize}
\item \textbf{Lemma 1:} if one \(A\)-step never leaves \(B\), then star also never leaves \(B\), and because star allows zero steps, star gives exactly \(B\).
\item \textbf{Lemma 2:} every star-result is either already in \(B\) because zero steps were taken, or it comes from making one genuinely new \(A\)-step and then continuing with star.
\end{itemize}

That is why the semantic lemmas mirror the operational definition of STAR so closely.

\end{document}
